\section{Limitations and Threats to Validity}
\label{sec:limitations}
The corpus has three synthetic families. Audits improve internal control, but
they do not show production realism. Incident 003 v3 was designed after an
earlier version exposed non-temporal shortcuts. Its audit covers known
shortcuts, not every learned cue.

Most comparisons use development data. Incident 001 has only raw validation.
Incident 002 has six constructions and a separate transport-recovery
sensitivity analysis. Incident 003 has four paired constructions and one call
per case. Its decline occurs only on twins. A concurrency category prior and
model variability remain possible. The study uses one GPT-5.6 model family.
It has no human baseline, matched validation comparison, or held-out result.

The taxonomy and component oracle are research choices. Component labels are
less stable in parts of Incidents 001 and 002. Evidence-support v5 was made
after Stage 1. Its Stage-1 use is post hoc; v4 remains the primary historical
rule.

The exposure vector measures representation, not practical attack success.
It does not test auxiliary data, membership inference, compromised HMAC keys,
provider retention, or legal compliance. Deterministic methods retain linkage
and frequency. Format preservation leaks shape. Relative time can be
reconstructed with an external anchor. Unannotated fields can remain visible.
Incorrect sensitivity labels can also expose data.

The repository does not publish provider JSONL files. It supports aggregate
reproduction, not exact response inspection. The study also excludes prompt
injection, malicious agents, poisoning, and insiders. Stronger claims require
more incidents, models, calls, human baselines, governance review, and a
preregistered held-out evaluation.
